\subsection{高频失分点}

本节不再逐篇搜集疏漏，也不虚构出现频数。以下八项来自 59 篇优秀论文中反复需要闭合的写作
输入输出关系，同时也是竞赛中一旦缺失便会削弱说服力的质量门。检查时只判断自己的稿件是否闭合。

\begin{enumerate}
  \item \textbf{题问没有闭合。}摘要、问题分析、模型、结果表和附件应逐问对应；每问至少留下
  “对象--方法--结果--解释”四个可核对位置。
  \item \textbf{选择没有来由。}写出模型或算法前，先说明观察到的结构、上一方法的具体不足、
  可行域规模或误差目标；删掉“综合考虑”“为了提高精度”后仍应看得懂为何换法。
  \item \textbf{数学对象没有落到题意。}变量、索引、单位、边界和约束都应能翻译回真实对象；
  公式后用一句话说明各项作用，结果后报告活跃约束或余量。
  \item \textbf{求解过程不可复现。}正文至少给初始化、参数、停止规则和输出；随机算法还要给
  预算、种子或重复统计，精确算法给界、间隙或可手算小例。
  \item \textbf{结果只有数字没有判断。}按“条件--读数--比较--原因--行动”解释图表；遇到跳变、
  异常时段或反直觉排序，先查事件对象、局部结构和统计定义变化，不用“结果合理”结尾。
  \item \textbf{检验与主模型同源。}机理题用守恒、极限、网格收敛或独立几何核对；优化题用
  约束回算、基线和共同情景复评；数据题按实体或时间隔离验证，并同时报告关键少数类。
  \item \textbf{不确定性没有传到决策。}参数扰动、预测概率、情景分布和中间分类误差应进入
  目标或约束；区分固定策略复算与重新优化，并给稳定区间或策略切换点。
  \item \textbf{正文与附件断开。}表图、代码和数据产物使用同一版本标识；正文给入口和关键
  参数，附录给依赖、数据字典、运行顺序与完整输出，修改结果后重新生成表图和摘要数字。
\end{enumerate}

这些门的用途是约束自己的稿件，不是给 59 篇论文排错。真正应模仿的仍是作者如何把观察写成
判断、把判断写成动作，再用结果和检核对住这一动作。
